\begin{lemma}
For every \(\delta>0\), and for every parameter tuple satisfying
\[
\noderef{ParameterHierarchy}(\varepsilon,\varepsilon_1,\varepsilon_2,n_0),
\]
there is a threshold \(N\ge n_0\) such that, for every \(n>N\), every
\(m,p\) with
\[
M(n)=\noderef{TwoBiteNaturalReducedVertexCount}(n),\qquad
m=M(n),\qquad
p=\noderef{TwoBiteEdgeProbability}\!\left(\frac12,n\right),
\]
and every set \(I\subseteq\mathrm{Fin}(n)\) with
\[
|I|=\noderef{TwoBiteNaturalIndependenceScale}(1+\varepsilon,n),
\]
one has
\[
\noderef{TwoBiteEventProbability}\!\left(n,m,p,\,
\{\omega:
I\text{ is independent in }\noderef{TwoBiteFinalGraph}(\omega)_0
\text{ and }
\noderef{TwoBiteRegularityEvent}
(|I|,\varepsilon,\varepsilon_1,\varepsilon_2,\omega)\}\right)
\le
\delta\binom n{|I|}^{-1}.
\]
\end{lemma}
\begin{proof}
Fix \(\delta>0\) and assume
\(\noderef{ParameterHierarchy}(\varepsilon,\varepsilon_1,\varepsilon_2,n_0)\).
Unfolding the hierarchy gives, in particular,
\[
0<\varepsilon_2<\varepsilon_1<\varepsilon<1
\]
and
\[
\noderef{ParameterHierarchyEventualComparisons}
(\varepsilon,\varepsilon_1,\varepsilon_2,n_0).
\]
Consequently
\[
0<\varepsilon,\qquad 0\le\varepsilon_1,\qquad \varepsilon_1\le1. \tag{H}
\]

Apply \(\noderef{FixedSetProjectionCellSummationBound}\) to the same
\(\delta\) and the present parameter hierarchy.  It gives an internal error
\(\eta>0\) and a threshold \(N_1\ge n_0\) such that the projection-cell double
sum is at most \(\delta\binom nK^{-1}\) for every \(n>N_1\).  Apply
\(\noderef{ProjectionSizeProbabilityBound}\) with this \(\eta\) and
\(\varepsilon\), using \(\eta>0\) and \(0<\varepsilon\); let \(N_2\) be its
threshold.  Define
\[
N=\max\{n_0,N_1,N_2\}.
\]
Then \(N\ge n_0\).  If \(n>N\), then
\[
n_0<n,\qquad N_1<n,\qquad N_2\le n, \tag{T}
\]
so all child thresholds are available.

Now fix \(n,m,p\) and \(I\subseteq\mathrm{Fin}(n)\) satisfying the hypotheses
of the theorem.  Put
\[
K=\noderef{TwoBiteNaturalIndependenceScale}(1+\varepsilon,n),
\qquad
p_0=\noderef{TwoBiteEdgeProbability}\!\left(\frac12,n\right).
\]
The active hypotheses give
\[
|I|=K,\qquad m=\noderef{TwoBiteNaturalReducedVertexCount}(n),\qquad p=p_0.
\tag{K}
\]
Let \(\mathcal B_I(\omega)\) be the event that \(I\) is independent in
\(\noderef{TwoBiteFinalGraph}(\omega)\).  Let
\[
\mathcal R_I(\omega)=
\noderef{TwoBiteRegularityEvent}(|I|,\varepsilon,\varepsilon_1,\varepsilon_2,
\omega),
\]
which is the regularity event appearing in the Lean statement.  By \(|I|=K\),
this event is the same as
\[
\mathcal R_K(\omega)=
\noderef{TwoBiteRegularityEvent}(K,\varepsilon,\varepsilon_1,\varepsilon_2,
\omega). \tag{R}
\]

For \(0\le\ell_R,\ell_B\le K\), define the projection-size cell
\[
\mathcal S_{\ell_R,\ell_B}(\omega)=
\noderef{TwoBiteProjectionSizeEvent}(I,K,\ell_R,\ell_B,\omega).
\]
Unfolding \(\noderef{TwoBiteProjectionSizeEvent}\), this cell records exactly
\[
|I|=K,\qquad
|\pi_R^\omega(I)|=\ell_R,\qquad
|\pi_B^\omega(I)|=\ell_B. \tag{S}
\]
Because \(|I|=K\), every sample has unique projection-image sizes
\(|\pi_R^\omega(I)|\) and \(|\pi_B^\omega(I)|\), each at most \(K\).  Hence
the cells \(\mathcal S_{\ell_R,\ell_B}\), for
\(\ell_R,\ell_B\in\{0,\ldots,K\}\), form a finite disjoint partition of the
sample space.  Since
\(\noderef{TwoBiteEventProbability}\) is the sum of sample weights over the
finite satisfying set, finite additivity over this disjoint partition gives
\[
\Pr(\mathcal B_I\cap\mathcal R_I)
=
\sum_{\ell_R=0}^{K}\sum_{\ell_B=0}^{K}
\Pr(\mathcal B_I\cap\mathcal R_K\cap\mathcal S_{\ell_R,\ell_B}),
\tag{1}
\]
where we used the rewrite (R) in the numerator event.

Fix a cell.  If \(\ell_R\ell_B<K\), then
\(\mathcal S_{\ell_R,\ell_B}\) is empty.  Indeed, the injective image
\(\pi_\omega(I)\) has \(K\) points, while (S) would place it inside the
rectangle \(\pi_R^\omega(I)\times\pi_B^\omega(I)\), which has only
\(\ell_R\ell_B<K\) points.  Therefore the corresponding summand in (1) is
\(0\), matching the zero branch in the summation helper.

Assume now \(\ell_R\ell_B\ge K\).  If
\(\Pr(\mathcal S_{\ell_R,\ell_B})=0\), then
\[
\Pr(\mathcal B_I\cap\mathcal R_K\cap\mathcal S_{\ell_R,\ell_B})=0
\]
because the numerator event is a subset of the cell.  Otherwise the definition
of \(\noderef{TwoBiteConditionalProbability}\) gives the exact factorization
\[
\Pr(\mathcal B_I\cap\mathcal R_K\cap\mathcal S_{\ell_R,\ell_B})
=
\Pr(\mathcal S_{\ell_R,\ell_B})\,
\Pr(\mathcal B_I\cap\mathcal R_K\mid\mathcal S_{\ell_R,\ell_B}). \tag{2}
\]
The conditional factor is at most \(1\).  It is also bounded by
\(\noderef{FixedSetConditionalIndependenceBound}\), applied with
\[
k=K,\qquad \ell_R=\ell_R,\qquad \ell_B=\ell_B,
\]
the same \(I,n,m,p,\varepsilon,\varepsilon_1,\varepsilon_2,n_0\), and the
hypotheses (H), (T), and (K).  The regularity event in the child numerator is
\(\mathcal R_K\), which is exactly the rewritten version of the active event
from (R).  Thus
\[
\Pr(\mathcal B_I\cap\mathcal R_K\mid\mathcal S_{\ell_R,\ell_B})
\le
\min\!\left\{1,\,
\exp\!\left(
-p\left(
\noderef{ProjectionOpenPairFunction}
((\ell_R:\mathbb R),(\ell_B:\mathbb R),(K:\mathbb R),p,(n:\mathbb R))
-\varepsilon^3(K:\mathbb R)^2\right)\right)\right\}. \tag{3}
\]
The minimum combines the trivial \(1\)-bound with the child exponential bound.

Next apply \(\noderef{ProjectionSizeProbabilityBound}\) to the cell.  The
threshold \(N_2\le n\), the rounded \(m\)-identity in (K), the equalities
\(|I|=K\) and
\[
K=\noderef{TwoBiteNaturalIndependenceScale}(1+\varepsilon,n),
\]
and the nonempty-cell hypothesis \(\ell_R\ell_B\ge K\) are precisely its
hypotheses.  Therefore, with
\[
x_R=\frac{\ell_R}{K},\qquad x_B=\frac{\ell_B}{K},
\]
we have
\[
\Pr(\mathcal S_{\ell_R,\ell_B})
\le
\exp\!\left(
\left(\eta-\frac{2-x_R-x_B}{2}\right)(K:\mathbb R)\log(n:\mathbb R)
\right). \tag{4}
\]

Combining (2), (3), and (4), every nonempty cell contributes at most
\[
\exp\!\left(
\left(\eta-\frac{2-(\ell_R:\mathbb R)/(K:\mathbb R)
        -(\ell_B:\mathbb R)/(K:\mathbb R)}{2}\right)
  (K:\mathbb R)\log(n:\mathbb R)\right)
\min\!\left\{1,\,
\exp\!\left(
-p\left(
\noderef{ProjectionOpenPairFunction}
((\ell_R:\mathbb R),(\ell_B:\mathbb R),(K:\mathbb R),p,(n:\mathbb R))
-\varepsilon^3(K:\mathbb R)^2\right)\right)\right\}.
\tag{5}
\]
Together with the empty-cell \(0\) case, substituting (5) into (1) gives
exactly the finite sum
\[
\sum_{\ell_R=0}^{K}\sum_{\ell_B=0}^{K}
\begin{cases}
0, & \ell_R\ell_B<K,\\
\text{the expression in (5)}, & \ell_R\ell_B\ge K.
\end{cases}
\tag{6}
\]
After rewriting \(p=p_0=\noderef{TwoBiteEdgeProbability}(\frac12,n)\), (6) is
the double sum appearing in
\(\noderef{FixedSetProjectionCellSummationBound}\), with its
\[
K=\noderef{TwoBiteNaturalIndependenceScale}(1+\varepsilon,n)
\]
and with all natural parameters cast to real numbers in the same positions:
the two projection sizes, \(K\), the \(n\)-argument of
\(\noderef{ProjectionOpenPairFunction}\), the square
\((K:\mathbb R)^2\), and the logarithmic factor.

Since \(N_1<n\), the summation helper now yields
\[
\Pr(\mathcal B_I\cap\mathcal R_I)
\le
\delta\,(\binom nK:\mathbb R)^{-1}. \tag{7}
\]
Finally \(K=|I|\) by (K), so
\[
(\binom nK:\mathbb R)^{-1}
=
\left(\binom n{|I|}:\mathbb R\right)^{-1}.
\]
Substituting this in (7) is exactly the bound in the Lean statement.
\end{proof}
